\documentclass[11pt,letterpaper]{article}
\usepackage[hmargin=1in,top=1in,bottom=1in]{geometry}
\usepackage{amsmath}
\usepackage{amssymb}
\usepackage{multirow}
\usepackage{libertinus}
\usepackage{natbib}

\usepackage{authblk}
\usepackage{booktabs}
\usepackage{dirtytalk}
\usepackage[hidelinks]{hyperref}
\usepackage{microtype}
\usepackage{multicol}
\usepackage{pifont}
\usepackage{simplex}
\usepackage{setspace} 
\usepackage{tikz}
\usepackage{url}

\bibpunct[: ]{(}{)}{;}{a}{}{,}
\usetikzlibrary{matrix}

\newcommand{\concat}{\ensuremath{\mathord{\smallfrown}}}
\newcommand{\gap}{\rule{1em}{0.5pt}}
\newcommand{\word}[1]{\textit{#1}}
\newcommand{\term}[1]{\textit{#1}}
\newcommand{\feature}[1]{\textsc{#1}}
\newcommand{\featspec}[2]{\ensuremath{#1}\textsc{#2}}
\newcommand{\y}{\ding{51}}
\newcommand{\n}{\ding{55}}
\newcommand*\circled[1]{\tikz[baseline=(char.base)]{
            \node[shape=circle,draw,inner sep=2pt] (char) {#1};}}
\newcommand{\writing}[1]{\ensuremath{\langle}#1\ensuremath{\rangle}}

% I am using this for anonymization.
\newif\ifanon
\anonfalse

% This name is sort of lame but I don't have anything clever yet.

\ifanon
\else
\title{Romanian palatalization without tears} 
\author{Kyle Gorman \& John Winstead}
\affil{CUNY Graduate Center, New York}
\date{}
\fi

\begin{document}
\pagestyle{empty}
\thispagestyle{empty}

\ifanon
\begin{center}
\textbf{
Romanian palatalization without tears
}
\end{center}
\else
\maketitle
\fi

In Romanian, coronal and dorsal obstruents become palatalized when followed by certain suffixes beginning with a front vowel, such as noun plurals in \word{-i}
(NB: /i/ $\rightsquigarrow$ [ʲ] word-finally).
\vspace{-\topsep}
% Get rid of z examples...
\begin{unlabeledexample}
\begin{tabular}{lll l lll l lll}
ru{s}   & ru{ʃ}ʲ   & `Russian(s)'  && bra{d}    & bra{z}ʲ   & `fir(s)'
                                    && mun{t}e & mun{ts}ʲ & `mountain(s)' \\
pro{s}t & pro{ʃ}tʲ & `fool(s)'     && miloɡ & milo{dʒ}ʲ & `beggar(s)'
                                    && pisi{k}ə  & pisi{tʃ}ʲ & `cat(s)' \\
\end{tabular}
\label{e:noun-plurals}
\end{unlabeledexample}

\noindent
These \term{palatalization} mutations are straightforwardly handled by a series of ordinary phonological rules or constraints.
However, additional observations have led some \citep[e.g.,][]{Chitoran:01} to posit that palatalization is morphophonological in part.
Palatalization exhibits \term{non-derived environment blocking} (NDEB):
it does not apply when the obstruent target and front vowel trigger are underlyingly adjacent and tautomorphemic;
e.g., [{ki}loɡram] `kilogram', [{ɡi}mpe] `thorn'.
% I am not sure if we want to consider these problematic for NDEB or not.
%However, there are a few nouns which end in /i/ in both the citation form
%(the nom.sg.~indef.)~and plural, and these also do not palatalize; e.g.,
%[privegʲ] `wake(s)', [triuŋgʲ] `triangle(s)', [uŋkʲ] `uncle(s)',
%[pəduke, pədukʲ] `louse--lice'.
But there are also some derived exceptions,
%While the exact analysis of these nouns is unclear, there are other
%derivational exceptions, 
such as some agent nominals in \word{-ist};
e.g., [fok-ist] `stoker', [taŋk-ist] `tank driver', [droɡ-ist] `druggist', [joɡ-ist] `yogi';
cf.~the palatalization observed in [kajak, kajatʃ-ist] `kayak(er)', [stɨŋɡ, stɨŋdʒ-ist] `left(ist)'.
\citet{Steriade:08} further claims that nouns that select plural allomorphs which trigger palatalization (e.g., \word{-i}) tend to select palatalizing allomorphs of derivational suffixes%
---like the homophonous \word{-i} verbalizer as in (\ref{e:steriade}a)---%
and nouns which select non-palatalizing plural allomorphs (e.g., \word{-uri}) tend to select non-palatalizing verbalizer allomorphs (e.g., \word{-a} or \word{-ui}) as in (\ref{e:steriade}b),
though there are exceptions to this generalization as in (\ref{e:steriade}c).

\begin{unlabeledexample}
\begin{tabular}{l lll ll}
a. & kola{k}  & kola{tʃ}ʲ & `kalach(es)' & ɨŋkola{tʃ}ʲ & `to coil' \\
   & verde    & verzʲ     & `green(s)'   & ɨnverzʲ     & `to become green' \\
b. & fok      & fokurʲ    & `fire(s)'    & ɨnfoka      & `to fire up' \\
   & tɨrɡ     & tɨrɡurʲ   & `market(s)'  & tɨrɡui      & `to go shopping' \\
c. & buluk    & bulukurʲ  & `crowd(s)'   & bulutʃʲ     & `to crowd (together)' \\
   & klopo{t} & klopo{t}e & `bell(s)'    & klopo{ts}ʲ  & `to ring a bell' \\
%- <oțet, oțeturi> `vinegar' but <oțeti> `to make into vinegar'
%- <scut, scuturi> `shield' but <scuti> `to exonerate'
%d. & pro{s}t  & pro{ʃ}tʲ  & `fool(s)'    & prostʲ      & `to fool' \\
\end{tabular}
\label{e:steriade}
\end{unlabeledexample}

We first consider the productivity of the different forms of palatalization.
Productivity analyses are conducted using a lexicon extracted from Wiktionary and the \emph{Dicționarul explicativ al limbii române} \citep[DEX;][]{DEX}
filtered by frequency \citep{LCC} to simulate realistic lexical knowledge \citep[e.g.,][]{Kodner:19}.
Below we give the type frequencies of noun plural palatalization ($n$: \# of stems eligible for palatalization,
$e$: \# of stems which are exceptions to palatalization) for the 1,000 most common nouns in our sample, 
with majority-rules \citep[14]{Nida:49} and Tolerance Principle \citep{Yang:05,Yang:16} calculations,
showing all coronal- and dorsal-obstruent-final stems when followed by heterosyllabic \word{i} or \word{e}.
Under both thresholds, there is productive palatalization of \word{c} and \word{g} before \word{i} and \word{e},
and of \word{t}, \word{d}, and \word{s} before \word{i}.

\begin{unlabeledexample}
\begin{tabular}{lrrll lrrll}
 & $n$ & $e$ & majority? & tolerated?
 & & $n$ & $e$ & majority? & tolerated? \\
\midrule
\ldots\word{c}-\word{i}\ldots  &  28 &   0 & \y & \y
 & \ldots\word{c}-\word{e}\ldots  &  14 &   0 & \y & \y \\
\ldots\word{g}-\word{i}\ldots  &  17 &   0 & \y & \y
 & \ldots\word{g}-\word{e}\ldots  &   3 &   0 & \y & \y \\
\ldots\word{t}-\word{i}\ldots  & 207 &   0 & \y & \y
 & \ldots\word{t}-\word{e}\ldots  & 229 & 229 & \n & \n \\
\ldots\word{d}-\word{i}\ldots  &  23 &   0 & \y & \y
 & \ldots\word{d}-\word{e}\ldots  &  36 &  36 & \n & \n \\
\ldots\word{s}-\word{i}\ldots  &  11 &   0 & \y & \y
 & \ldots\word{s}-\word{e}\ldots  &  44 &  44 & \n & \n \\
\ldots\word{z}-\word{i}\ldots  &   2 &   1 & (n.a.) & (n.a.)
 & \ldots\word{z}-\word{e}\ldots  &  48 &  48 & \n & \n \\
\ldots\word{st}-\word{i}\ldots &   4 &   0 & \y & \y
 & \ldots\word{st}-\word{e}\ldots &   0 &   0 & (n.a.) & (n.a.) \\
\ldots\word{ct}-\word{i}\ldots &   3 &   0 & \y & \y
 & \ldots\word{ct}-\word{e}\ldots &  22 &  22 & \n & \n \\
\end{tabular}
\label{e:productivity}
\end{unlabeledexample}

We argue that the \say{morphological} effect observed by \citet{Steriade:08} can, instead, be derived within the narrow phonology without reference to morphological structure.
Both NDEB effects and apparent exceptions like the \word{-ist} class stem from coronal and dorsal obstruents which are prespecified and therefore inalterable,
whereas palatalizing obstruents are underspecified---echoing the notion of \term{inalterability as prespecification}
\citep[e.g.,][]{Inkelas:Cho:93,Kiparsky:93,Buckley:94,Inkelas:95,Inkelas:EtAl:97,Inkelas:00,Inkelas:14,Gorman:26}.
\citet{Rasin:25} provides an analysis of Romanian NDEB in this vein:
segments which undergo palatalization are underspecified and thus mutable w.r.t.~feature-filling palatalization rules,
but prespecified segments are inalterable w.r.t.~these rules.
To implement the desired feature-filling behavior, we use unification rules \citep{Bale:EtAl:14},
formalized in Logical Phonology
\citep[e.g.,][]{Bale:Reiss:18,Bale:EtAl:20,Reiss:21,Dabbous:EtAl:24,Leduc:EtAl:24,Reiss:24,Gorman:Reiss:25:Metaphony}.

Specifically, the dorsals /K/ and /G/ are underspecified for \feature{Coronal} and palatalize before any \featspec{+}{front} vowel,
while /k/ and /g/ are prespecified \featspec{-}{Coronal} and therefore inalterable.
Rules (\ref{e:velar-pal}--\ref{e:default}) map underspecified /K, G/ to [tʃ, dʒ] (resp.)~before front vowels and to [k, ɡ] elsewhere.
While these rules intensionally target many segments,
unification is defined so that it does not overwrite existing feature specifications,
so prespecified segments like /k, ɡ/ are inalterable w.r.t.~both rules.
The coronals /T/, /D/, and /S/ palatalize only before /i/, with complementary underspecification:
/T/ and /D/ lack \feature{strident} and \feature{Continuant},
while /S/ lacks \feature{anterior} and \feature{distributed}.
Fully specified /z/ is inalterable in all contexts.

The tendency for plural and verbalizer allomorph selection to co-vary follows directly from the same feature specifications that drive inalterability.
The non-palatalizing allomorphs (\textit{-uri}, \textit{-a}, \textit{-ui}) select for roots specified
[\featspec{-}{Coronal}, \featspec{+}{Dorsal}, \featspec{-}{Continuant}]---the very specification that renders prespecified obstruents inalterable to palatalization---while the palatalizing allomorphs (\textit{-i}, \textit{-e}) are defaults bled by more specific---in the Pāṇinian sense---phonological rules.
Since spell-out is inside-out \citep{Bobaljik:00},
\textsc{Num} and \textit{v} are spelled out after the root,
giving them access to its underlying form.
The phonologically null \textit{n} head in
[\textsc{NumP} [\textit{n}P $\surd$\textsc{root} \textit{n}] \textsc{Num}[\textsc{pl}]]
and [\textit{v}P [\textit{n}P $\surd$\textsc{root} \textit{n}] \textit{v}] is \say{pruned} \citep{Embick:03} and therefore transparent,
so no intervening head blocks this visibility.
The co-variation \citeauthor{Steriade:08} attributes to paradigm consultation is thus epiphenomenal:
both the plural allomorph and the verbalizer allomorph are conditioned by the same underlying root specification, with no reference to inflectional identity required.
The formal spell-out rules and a sample derivation are given below.

\noindent
\begin{minipage}[t]{.6\textwidth}
\begin{example}[Dorsal palatalization]
$\begin{bmatrix} 
\featspec{-}{Syllabic} \\ 
\featspec{+}{Consonantal} \\
\featspec{-}{Sonorant} \\ 
\featspec{-}{Approximant} \\
\featspec{-}{Nasal} \\ 
\ldots 
\end{bmatrix}~\sqcup~\begin{Bmatrix}
\featspec{+}{Coronal} \\ 
\featspec{-}{Dorsal} \\
\featspec{-}{Anterior} \\ 
\featspec{+}{Distributed} \\
\featspec{+}{Strident} 
\end{Bmatrix}~\gap~/~\begin{bmatrix}
\featspec{+}{Syllabic} \\ 
\featspec{+}{Front} 
\end{bmatrix}$
\label{e:velar-pal}
\end{example}
\end{minipage}%
~%
\begin{minipage}[t]{.35\textwidth}
\begin{example}[Dorsal default fill-in]
$\begin{bmatrix}
\featspec{-}{Syllabic} \\ 
\featspec{+}{Consonantal} \\
\featspec{-}{Sonorant} \\ 
\featspec{-}{Approximant} \\
\featspec{-}{Nasal} \\ 
\ldots 
\end{bmatrix}~\sqcup~\begin{Bmatrix}
\featspec{-}{Coronal} \\ 
\featspec{+}{Dorsal} \\
\featspec{-}{Continuant} 
\end{Bmatrix}$
\label{e:default}
\end{example}
\end{minipage}
\vspace{\baselineskip}

\begin{multicols}{2}
\begin{unlabeledexample}
$\text{Num}[\feature{Pl}]~\Longleftrightarrow~
\begin{cases}
\textrm{/-uri/} & \begin{bmatrix}
    \featspec{-}{Coronal} \\
    \featspec{+}{Dorsal} \\
    \featspec{-}{Continuant}
\end{bmatrix}~\gap \\[3pt]
\textrm{/-i/} & \textrm{elsewhere}
\end{cases}$
\label{e:plural-spellout}
\end{unlabeledexample}

\begin{unlabeledexample}
$\word{v}~\Longleftrightarrow~
\begin{cases}
\textrm{/-a/} & \begin{bmatrix}
    \featspec{-}{Coronal} \\
    \featspec{+}{Dorsal} \\
    \featspec{-}{Continuant}
\end{bmatrix}~\gap \\[3pt]
\textrm{/-i/} & \textrm{elsewhere}
\end{cases}$
\label{e:velar-spellout}
\end{unlabeledexample}
\end{multicols}

\noindent
The interaction between allomorph selection and palatalization is illustrated in the derivations below,
which begin immediately after root spell-out.

\begin{unlabeledexample}
\begin{tabular}{l llllll}
UR: & fok & fok$\concat\feature{Num}[\feature{Pl}]$
    & koloK & koloK$\concat\feature{Num}[\feature{Pl}]$
    & miloG & miloG$\concat v$ \\
\midrule
Rule (\ref{e:plural-spellout}): &     & fokuri &       & koloKi  &       &         \\
Rule (\ref{e:velar-spellout}):  &     &        &       &         &       & miloGi  \\
Rule (\ref{e:velar-pal}):       &     &        &       & kolatʃi &       & milodʒi \\
Rule (\ref{e:default}):         &     &        & kolak &         & miloɡ &         \\
Desyllabification:              &     & fokurʲ &       & kolatʃʲ &       & milodʒʲ \\
\midrule
SR: & fok & fokurʲ & kolak & kolatʃʲ & miloɡ & milodʒʲ \\
\end{tabular}
\label{e:derivation}
\end{unlabeledexample}

\newpage

\bigskip\noindent
The coronals show the same split as the dorsals, two mutations before \word{-i} that apply to some segments and skip others, and the skipping again follows from the root.
The prespecified segments carry the value a mutation would supply and so resist it, while their underspecified counterparts, lacking that value, undergo.
The specifications are given in (\ref{e:cor-inventory}); /T, D/ and /S, Z/ are the underspecified stops and fricatives, /t, d/ and /s, z/ the prespecified ones.

\begin{unlabeledexample}
\small
\begin{tabular}{l cccccc}
          & /t, d/ & /T/ & /D/ & /s, z/ & /S, Z/ \\
\midrule
\feature{Coronal}    & $+$ & $+$ & $+$ & $+$ & $+$ \\
\feature{Anterior}   & $+$ & $+$ & $+$ & $+$ &     \\
\feature{Strident}   & $-$ &     &     & $+$ & $+$ \\
\feature{Continuant} & $-$ & $-$ &     & $+$ & $+$ \\
\end{tabular}
\label{e:cor-inventory}
\end{unlabeledexample}

\noindent
The one asymmetry to note now is that /T/ is prespecified \featspec{-}{Continuant} while /D/ is open for \feature{Continuant}, which later splits [ts] from [z].

\bigskip\noindent
Before \word{-i}, the strident fricatives shift to postalveolar, /s, z/ surfacing as [ʃ, ʒ].
\term{S-palatalization} (\ref{e:s-pal}) unifies postalveolar place into the coronal stridents, so underspecified /S, Z/ take it and prespecified /s, z/, already \featspec{+}{Anterior}, resist.
The target locates /i/ by a rightward SEARCH terminating on the first [\featspec{+}{High}, \featspec{+}{Front}] segment \citep{Dabbous:EtAl:21,Dabbous:EtAl:24,Azadegan:26}.
Since only the vowel terminates the scan, an intervening stop is passed over, and /S/ palatalizes across it in \word{prost}.

\begin{unlabeledexample}
\small\hfill$\displaystyle
\begin{bmatrix}
\featspec{+}{Coronal}\\
\featspec{+}{Strident}\\
\featspec{+}{Continuant}
\end{bmatrix}
~\sqcup~
\begin{Bmatrix}
\featspec{-}{Anterior}\\
\featspec{+}{Distributed}
\end{Bmatrix}
\;/\;\gap~
\begin{bmatrix}
\featspec{+}{High}\\
\featspec{+}{Front}
\end{bmatrix}
$\hfill\null
\label{e:s-pal}
\end{unlabeledexample}

\noindent
In an \word{st} cluster the stop is protected, so \word{prost} gives [proʃt] and not *[proʃts].
A bleeding rule (\ref{e:bleed}) delinks \feature{Strident} from the stop after the derived postalveolar, so the following assibilation finds nothing to strengthen.
The protection must be its own rule, since assibilation searches rightward to its trigger and cannot see the postalveolar lying to the stop's left.
The environment uses circle notation, \circled{S} denoting the natural class whose bundle is a superset of the derived postalveolar's, located here by a leftward SEARCH.

\begin{unlabeledexample}
\small\hfill$\displaystyle
\begin{bmatrix}
\featspec{+}{Coronal}
\end{bmatrix}
~\setminus~
\begin{Bmatrix}
\featspec{+}{Strident}
\end{Bmatrix}
\;/\;
\begin{bmatrix}
\circled{S}
\end{bmatrix}
~\gap
$\hfill\null
\label{e:bleed}
\end{unlabeledexample}

\noindent
Before \word{-i}, the coronal stops also assibilate to affricates.
\term{Assibilation} (\ref{e:assib}) unifies \featspec{+}{Strident} into an anterior coronal stop, so underspecified /T, D/ assibilate and prespecified /t, d/, already \featspec{-}{Strident}, resist, as does any stop the bleed (\ref{e:bleed}) has already cleared.

\begin{unlabeledexample}
\small\hfill$\displaystyle
\begin{bmatrix}
\featspec{+}{Coronal}\\
\featspec{+}{Anterior}
\end{bmatrix}
~\sqcup~
\begin{Bmatrix}
\featspec{+}{Strident}
\end{Bmatrix}
\;/\;\gap~
\begin{bmatrix}
\featspec{+}{High}\\
\featspec{+}{Front}
\end{bmatrix}
$\hfill\null
\label{e:assib}
\end{unlabeledexample}

\noindent
Late default fill-ins (\ref{e:cor-default}), three context-free rules stated together for space, supply the unmarked values to whatever the mutations left open.
Unmutated stops surface as [t, d] and unmutated fricatives as [s, z].
The third rule fills \featspec{+}{Continuant} on assibilated stops, and here the /T/--/D/ asymmetry surfaces, /D/ spirantizes to [z] because it has no \feature{Continuant} value to block the fill, while prespecified /T/ stays [ts].
Romanian has no voiced coronal affricate, a gap that follows directly, since the only segment that could surface as [dz] is exactly the one the default spirantizes.

\begin{unlabeledexample}
\small\hfill$\displaystyle
\begin{bmatrix}
\featspec{+}{Coronal}
\end{bmatrix}
~\sqcup~
\begin{Bmatrix}
\featspec{-}{Strident}
\end{Bmatrix},~
\begin{bmatrix}
\featspec{+}{Coronal}
\end{bmatrix}
~\sqcup~
\begin{Bmatrix}
\featspec{+}{Anterior}\\
\featspec{-}{Distributed}
\end{Bmatrix},~
\begin{bmatrix}
\featspec{+}{Coronal}\\
\featspec{+}{Strident}
\end{bmatrix}
~\sqcup~
\begin{Bmatrix}
\featspec{+}{Continuant}
\end{Bmatrix}
$\hfill\null
\label{e:cor-default}
\end{unlabeledexample}

\noindent
A final rule (\ref{e:glide}) desyllabifies a word-final /i/ after a consonant, stripping \featspec{+}{Syllabic} so it surfaces as secondary palatalization on the preceding consonant.
It applies after the mutations, which are triggered by /i/ as a vowel, so the glide plays no role in conditioning them.

\begin{unlabeledexample}
\small\hfill$\displaystyle
\text{i}~\setminus~
\begin{Bmatrix}
\featspec{+}{Syllabic}
\end{Bmatrix}
~\sqcup~
\begin{Bmatrix}
\featspec{-}{Syllabic}
\end{Bmatrix}
\;/\;
\begin{bmatrix}
\featspec{+}{Consonantal}
\end{bmatrix}
~\gap~\#
$\hfill\null
\label{e:glide}
\end{unlabeledexample}

\bigskip\noindent
The same contrast conditions the spell-out of the plural and verbalizer suffixes, just as it did for the dorsals.
Let \textsc{p} abbreviate the prespecified obstruent context, the disjunction of the three inalterability conditions, one per mutation:
\[
\textsc{p} \;=\;
\begin{bmatrix}
\featspec{-}{Coronal}\\
\featspec{+}{Dorsal}\\
\featspec{-}{Continuant}
\end{bmatrix}
\;\textrm{or}\;
\begin{bmatrix}
\featspec{+}{Coronal}\\
\featspec{-}{Strident}
\end{bmatrix}
\;\textrm{or}\;
\begin{bmatrix}
\featspec{+}{Coronal}\\
\featspec{+}{Anterior}\\
\featspec{+}{Strident}
\end{bmatrix}
\]

\begin{multicols}{2}
\begin{unlabeledexample}
$\text{Num}[\feature{Pl}]~
\Longleftrightarrow~
\begin{cases}
\textrm{/-uri/} & /~
\textsc{p} \\[3pt]
\textrm{/-i/} & 
\textrm{elsewhere}
\end{cases}$
\label{e:plural-spellout}
\end{unlabeledexample}

\begin{unlabeledexample}
$\word{v}~\Longleftrightarrow~
\begin{cases}
\textrm{/-a/} & /~\textsc{p} \\[3pt]
\textrm{/-i/} & \textrm{elsewhere}
\end{cases}$
\label{e:verb-spellout}
\end{unlabeledexample}
\end{multicols}

\noindent
The dorsal disjunct is the one from the earlier spell-out; the two coronal disjuncts are the inalterability conditions of assibilation (\ref{e:assib}) and S-palatalization (\ref{e:s-pal}).
The exponents thus select the non-palatalizing allomorph for exactly the roots the phonology leaves unpalatalized, now at every place of articulation.

\begin{unlabeledexample}
\begin{tabular}{l lllll}
UR: & ruS-\feature{Num}[\feature{Pl}]
    & proST-\feature{Num}[\feature{Pl}]
    & braD-\feature{Num}[\feature{Pl}]
    & skut-\feature{Num}[\feature{Pl}]
    & klopoT-\word{v} \\
\midrule
Rule (\ref{e:plural-spellout}):  & ruSi  & proSTi  & braDi  & skuturi  &          \\
Rule (\ref{e:verb-spellout}):    &       &         &        &          & klopoTi  \\
Rule (\ref{e:s-pal}):            & ruʃi  & proʃTi  &        &          &          \\
Rule (\ref{e:bleed}):            &       & proʃTi  &        &          &          \\
Rule (\ref{e:assib}):            &       &         & braDzi &          & klopotsi \\
Rule (\ref{e:cor-default}):      &       &         & brazi  &          &          \\
Rule (\ref{e:glide}):            & ruʃj  & proʃtj  & brazj  & skuturj  & klopotsj \\
\midrule
SR: & ruʃʲ & proʃtʲ & brazʲ & skuturʲ & klopotsʲ \\
\end{tabular}
\label{e:coronal-derivation}
\end{unlabeledexample}

%New stuff I'm working through. Not super relevant to the prespecification analysis proper.

\bigskip\noindent
The rules so far derive palatalization before \word{-i}.
Two further facts push the same rules to more general forms, the categorical \word{sc} alternation before \word{-e} and the non-palatalizing \word{ct} loans.

\bigskip\noindent
\word{sc}-final stems palatalize categorically before \word{-e}, and the whole cluster mutates, /sk/ surfacing as [ʃt] (e.g., [muskə, muʃte] `fly(ies)', [broaskə, broaʃte] `frog(s)').
Two things about this are strange.
First, it looks like metathesis, since the places seem reordered, [ʃ] before [t] where the underlying cluster has /s/ before /k/.
Second, an expected segment disappears.
Palatalizing the /s/ gives [ʃ] and palatalizing the /k/ before \word{-e} gives the affricate [tʃ], so the expected output is *[muʃtʃe], with two postalveolars, yet only one surfaces.
The account reads the cluster left to right.
The /s/ is not adjacent to the vowel, so it cannot take place from /e/ directly; the /k/ between them is adjacent, and the vowel palatalizes it to [tʃ].
That derived [tʃ] is then what gives the /s/ its postalveolar place, yielding [ʃ].
The bleeding rule clears the [tʃ] back to a plain stop, so the cluster surfaces as [ʃt] with no affricate left.
The affricate exists only to carry palatal place leftward from the /k/ to the /s/, then is erased, which is why the missing postalveolar and the apparent metathesis are one fact seen twice.

\bigskip\noindent
This derivation locates the S-palatalization trigger in place, not height.
The basic rule fires on a [\featspec{+}{High}, \featspec{+}{Front}] vowel, but neither /e/ nor the [tʃ] derived from /k/ is \featspec{+}{High}, so that rule would leave the /s/ untouched.
What the /s/ reaches for is the postalveolar place of the [tʃ], so the revised trigger is the postalveolar class [\featspec{+}{Coronal}, \featspec{-}{Anterior}, \featspec{+}{Distributed}] (\ref{e:s-pal-rev}), and the /s/ searches rightward to the first segment bearing it.

\begin{unlabeledexample}
\small\hfill$\displaystyle
\begin{bmatrix}
\featspec{+}{Coronal}\\
\featspec{+}{Strident}\\
\featspec{+}{Continuant}
\end{bmatrix}
~\sqcup~
\begin{Bmatrix}
\featspec{-}{Anterior}\\
\featspec{+}{Distributed}
\end{Bmatrix}
\;/\;\gap~
\begin{bmatrix}
\featspec{+}{Coronal}\\
\featspec{-}{Anterior}\\
\featspec{+}{Distributed}
\end{bmatrix}
$\hfill\null
\label{e:s-pal-rev}
\end{unlabeledexample}

\noindent
One trigger is meant to cover both environments, since the plain \word{-i} forms and the \word{sc}-\word{e} forms present the /s/ with the same postalveolar place, only on different segments.
In the plain forms the terminator is /i/ itself, a high front vocoid and so coronal, which bears the place inherently and halts the scan while still a vowel; desyllabification (\ref{e:glide}) applies later, as before.
In the \word{sc} forms the terminator is instead the derived [tʃ], so the derivation runs in an order, the dorsal rule palatalizes /k/ first and thereby feeds S-palatalization the postalveolar it searches for.
The vowel /e/ never terminates the scan on its own, being \featspec{+}{Anterior} and so outside the place class, which is why the /s/ in \word{muskə} takes its place from the [tʃ] to its right and not from the /e/ two segments over.
The asymmetry the basic rule drew from vowel height thus reduces to place, and /i/ was arguably a place trigger from the start, its height incidental to a front vocoid already being coronal.

\bigskip\noindent
The feeding relies on the dorsal rule firing in \word{muskə}, where the segment after /k/ is the vowel /e/.
That the same rule does \emph{not} fire in the \word{ct} loans, where a consonant follows the dorsal, is the subject of the next revision, and the two cases turn on that one difference in what sits to the dorsal's right.

\bigskip\noindent
The \word{sc} case also stretches the bleed rule.
Subtracting \feature{Strident} alone protects the \word{st} stop, which is already plain in place, but it cannot clean up \word{sc}, where the segment to clear is the derived [tʃ] with a full postalveolar specification, so \word{muskə} would still surface as *[muʃtʃe].
The bleed must strip the whole postalveolar place, \feature{Anterior}, \feature{Distributed}, and \feature{Strident}, and confine itself to consonants so it leaves the suffix vowel alone (\ref{e:bleed-rev}); the coronal default then refills the plain stop.

\begin{unlabeledexample}
\small\hfill$\displaystyle
\begin{bmatrix}
\featspec{+}{Coronal}\\
\featspec{+}{Consonantal}
\end{bmatrix}
~\setminus~
\begin{Bmatrix}
\text{Anterior}\\
\text{Distributed}\\
\text{Strident}
\end{Bmatrix}
\;/\;
\begin{bmatrix}
\circled{S}
\end{bmatrix}
~\gap
$\hfill\null
\label{e:bleed-rev}
\end{unlabeledexample}

\noindent
The subtrahend removes each feature at whatever value it carries, and \circled{S} is the postalveolar environment of the basic bleed.
One rule now handles both clusters, clearing the derived [tʃ] in \word{sc} and stripping the /T/'s \featspec{+}{Anterior} in \word{st} so assibilation can no longer reach it; either way the default supplies [t].

\bigskip\noindent
The \word{ct} loans bear on the dorsal rule instead.
Latinate stems in \word{ct} leave the dorsal unpalatalized before \word{-e} (e.g., [punkt, punkte] `point(s)'), yet a plain search for a front vowel would wrongly palatalize /K/ across the intervening /T/.
The block falls out once the dorsal rule carries a broad terminator, halting on the very next segment, with the bracket in (\ref{e:velar-pal}) a condition on that segment rather than a terminator class.
In \word{musk-e} the segment after /K/ is /e/, which meets [\featspec{+}{Syllabic}, \featspec{+}{Front}], so /K/ palatalizes; in \word{punkt-e} it is /T/, which halts the scan and fails the condition, so /K/ is blocked.
This is the terminator-versus-condition distinction of \citet{Dabbous:EtAl:21,Dabbous:EtAl:24}, a narrow terminator letting interveners through, a broad terminator with a condition making them block.

\begin{unlabeledexample}
\begin{tabular}{l ll}
UR:                         & muSKe   & punKTe \\
\midrule
Rule (\ref{e:velar-pal}):   & mustʃe  & punkte \\
Rule (\ref{e:s-pal-rev}):   & muʃtʃe  &        \\
Rule (\ref{e:bleed-rev}):   & muʃTe   &        \\
Rule (\ref{e:cor-default}): & muʃte   &        \\
\midrule
SR:                         & muʃte   & punkte \\
\end{tabular}
\label{e:cluster-derivation}
\end{unlabeledexample}

\bigskip\noindent
The two palatalizations locate their triggers the same way, by SEARCH, and differ in one setting, how broad the terminator is.
The dorsal rule stops at the next segment, so any intervening consonant blocks.
S-palatalization stops only at a postalveolar, so a non-postalveolar in between is not a stopping point at all.
That single difference is the whole of locality here, and it means S-palatalization needs no locality setting of its own, applying adjacently in \word{rus} and across the stop in \word{prost}, local on one stem and long-distance on another, with the difference lying entirely in what intervenes.